\section{Appendix: Key IR Walkthroughs}

The following is a composite function before and after the \texttt{my-default}
pipeline, drawn from the framework's own tests. It exercises the whole chain:
canonicalization aligns the two adds, simplification folds the zero adds and
strength reduces the multiplies, common subexpression elimination merges the now
identical shifts, and dead code elimination clears the dead xor.

\subsection{Before}

\begin{lstlisting}[style=ir]
define i32 @composite(i32 %x, i32 %y) {
  %a = add i32 %x, 0
  %b = add i32 0, %x
  %c = mul i32 %a, 4
  %d = mul i32 %b, 4
  %e = add i32 %c, %d
  %dead = xor i32 %y, %y
  ret i32 %e
}
\end{lstlisting}

\subsection{After}

\begin{lstlisting}[style=ir]
define i32 @composite(i32 %x, i32 %y) {
  %1 = shl i32 %x, 2
  %e = add i32 %1, %1
  ret i32 %e
}
\end{lstlisting}

The result computes the same value, eight times the input, with the redundancy,
the identities, and the dead computation all gone.

\subsection{Strength reduction with flag preservation}

A multiply by a constant power of two becomes a shift, and the no signed wrap
flag is carried onto the shift, which overflows exactly when the multiply did.

\begin{lstlisting}[style=ir]
; before
%r = mul nsw i32 %x, 16
; after
%r = shl nsw i32 %x, 4
\end{lstlisting}

\subsection{Reproducing the numbers}

Every figure and table in this report is generated by
\texttt{scripts/generate\_plots.py} from the committed
\texttt{benchmarks/results/benchmarks.csv} and \texttt{opcodes.csv}, which are
themselves produced by \texttt{scripts/run\_benchmarks.sh}. The report contains
no hand typed measurement.
